\textbf{Phần 1. Trắc nghiệm lựa chọn}
\setcounterpageref{ex}{0}

\begin{ex} %[1D2V07]
Hình tam giác có bao nhiêu cạnh?
\choice
{$2$}
{$4$}
{\True $3$}
{$1$}
\loigiai{
Hình tam giác là hình khép kín có 3 cạnh và 3 góc (đỉnh).
\textbf{Đáp số:} $3$.\par
\textbf{Đáp án của câu hỏi là C}
}
\end{ex}

\begin{ex} %[1D2V07]
Hình tam giác có bao nhiêu góc (đỉnh)?
\choice
{\True $3$}
{$2$}
{$4$}
{$1$}
\loigiai{
Mỗi hình tam giác có đúng 3 đỉnh.
\textbf{Đáp số:} $3$.\par
\textbf{Đáp án của câu hỏi là A}
}
\end{ex}

\begin{ex} %[1D2V07]
Hình vuông có bao nhiêu cạnh bằng nhau?
\choice
{$3$}
{$5$}
{$2$}
{\True $4$}
\loigiai{
Hình vuông có 4 cạnh bằng nhau và 4 góc vuông.
\textbf{Đáp số:} $4$.\par
\textbf{Đáp án của câu hỏi là D}
}
\end{ex}

\begin{ex} %[1D2V07]
Hình chữ nhật có bao nhiêu góc vuông?
\choice
{$3$}
{\True $4$}
{$5$}
{$2$}
\loigiai{
Hình chữ nhật có 4 góc vuông và 4 cạnh (hai cạnh dài bằng nhau, hai cạnh ngắn bằng nhau).
\textbf{Đáp số:} $4$.\par
\textbf{Đáp án của câu hỏi là B}
}
\end{ex}

\begin{ex} %[1D2V07]
Hình tròn có bao nhiêu đỉnh (góc nhọn)?
\choice
{$1$}
{$2$}
{\True $0$}
{$3$}
\loigiai{
Hình tròn được bao quanh bởi một đường cong khép kín trơn cong đều, không có cạnh thẳng và không có đỉnh (góc).
\textbf{Đáp số:} $0$.\par
\textbf{Đáp án của câu hỏi là C}
}
\end{ex}

\begin{ex} %[1D2V07]
Một hình vuông được kẻ một đường thẳng nối hai đỉnh đối diện (đường chéo). Hỏi hình vuông đó được chia thành bao nhiêu hình tam giác?
\choice
{\True $2$}
{$1$}
{$3$}
{$4$}
\loigiai{
Đường chéo chia hình vuông thành $2$ hình tam giác bằng nhau.
\textbf{Đáp số:} $2$.\par
\textbf{Đáp án của câu hỏi là A}
}
\end{ex}

\begin{ex} %[1D2V07]
Một hình chữ nhật được chia đôi bởi một đường thẳng nối trung điểm hai cạnh đối diện. Hỏi tạo thành bao nhiêu hình chữ nhật nhỏ?
\choice
{$1$}
{$3$}
{$4$}
{\True $2$}
\loigiai{
Đường thẳng chia hình chữ nhật lớn thành $2$ hình chữ nhật nhỏ hơn.
\textbf{Đáp số:} $2$.\par
\textbf{Đáp án của câu hỏi là D}
}
\end{ex}

\begin{ex} %[1D2V07]
Cần ít nhất bao nhiêu que tính bằng nhau để xếp thành một hình tam giác?
\choice
{$2$}
{\True $3$}
{$4$}
{$1$}
\loigiai{
Mỗi cạnh của hình tam giác cần 1 que tính, vậy cần ít nhất $3$ que tính.
\textbf{Đáp số:} $3$.\par
\textbf{Đáp án của câu hỏi là B}
}
\end{ex}

\begin{ex} %[1D2V07]
Cần ít nhất bao nhiêu que tính bằng nhau để xếp thành một hình vuông?
\choice
{$3$}
{$5$}
{\True $4$}
{$2$}
\loigiai{
Mỗi cạnh của hình vuông cần 1 que tính bằng nhau, vậy cần ít nhất $4$ que tính.
\textbf{Đáp số:} $4$.\par
\textbf{Đáp án của câu hỏi là C}
}
\end{ex}

\begin{ex} %[1D2V07]
Để xếp được hai hình tam giác rời nhau, em cần ít nhất bao nhiêu que tính?
\choice
{\True $6$}
{$5$}
{$7$}
{$4$}
\loigiai{
Mỗi hình tam giác rời nhau cần 3 que tính. Hai hình cần: $3 + 3 = 6$ (que tính).
\textbf{Đáp số:} $6$.\par
\textbf{Đáp án của câu hỏi là A}
}
\end{ex}

\begin{ex} %[1D2V07]
Để xếp hai hình tam giác có chung một cạnh, em cần ít nhất bao nhiêu que tính?
\choice
{$4$}
{$6$}
{$3$}
{\True $5$}
\loigiai{
Hình tam giác thứ nhất cần 3 que tính. Hình tam giác thứ hai dùng chung 1 cạnh nên chỉ cần thêm 2 que tính nữa: $3 + 2 = 5$ (que tính).
\textbf{Đáp số:} $5$.\par
\textbf{Đáp án của câu hỏi là D}
}
\end{ex}

\begin{ex} %[1D2V07]
Một tờ giấy hình chữ nhật có 4 góc. Dùng kéo cắt đi một góc theo đường chéo nhỏ. Hỏi mảnh giấy lớn còn lại có bao nhiêu góc?
\choice
{$4$}
{\True $5$}
{$6$}
{$3$}
\loigiai{
Ban đầu có 4 góc. Khi cắt mất 1 góc, chỗ cắt tạo thêm 2 góc mới, vậy số góc còn lại là: $4 - 1 + 2 = 5$ (góc).
\textbf{Đáp số:} $5$.\par
\textbf{Đáp án của câu hỏi là B}
}
\end{ex}

\textbf{Phần 2. Trắc nghiệm đúng sai}
\setcounterpageref{ex}{0}

\begin{ex} %[1D2V07]
Xét tính đúng hoặc sai của các khẳng định toán học sau đây:
\choiceTF
{\True Trong một bức tranh có: 3 hình vuông, 2 hình tròn và 4 hình tam giác. trong tranh có tất cả $9$ hình phẳng.}
{Có 5 hình vuông màu đỏ và 4 hình vuông màu xanh. có tất cả $8$ hình vuông.}
{\True Một chiếc bánh sinh nhật hình tròn được cắt làm đôi bởi 1 nhát dao thẳng qua tâm. tạo thành $2$ miếng bánh.}
{Một chiếc bánh sinh nhật hình tròn được cắt bởi 2 nhát dao thẳng vuông góc nhau qua tâm. chiếc bánh được chia thành $5$ miếng.}
\loigiai{
\begin{itemize}
    \item Ý a) Đúng: Tổng số hình phẳng có trong bức tranh là: $3 + 2 + 4 = 9$ (hình).
\textbf{Đáp số:} $9$.
    \item Ý b) Sai: Kết quả đúng là $9$, không phải $8$. Tổng số hình vuông là: $5 + 4 = 9$ (hình vuông).
\textbf{Đáp số:} $9$.
    \item Ý c) Đúng: Một nhát dao thẳng cắt qua tâm chia hình tròn thành $2$ nửa hình tròn bằng nhau.
\textbf{Đáp số:} $2$.
    \item Ý d) Sai: Kết quả đúng là $4$, không phải $5$. Hai nhát dao cắt vuông góc qua tâm chia hình tròn thành $4$ phần bằng nhau.
\textbf{Đáp số:} $4$.
\end{itemize}
\begin{itemize}
    \item a) \textbf{Đúng}.
    \item b) \textbf{Sai}.
    \item c) \textbf{Đúng}.
    \item d) \textbf{Sai}.
\end{itemize}
}
\end{ex}

\begin{ex} %[1D2V07]
Xét tính đúng hoặc sai của các bài toán và phép tính sau đây:
\choiceTF
{Có 10 que tính bằng nhau, bạn Hùng xếp được 2 hình vuông riêng biệt. Hùng còn thừa $3$ que tính.}
{\True Có 10 que tính bằng nhau, bạn Mai xếp được 3 hình tam giác riêng biệt. Mai còn thừa $1$ que tính.}
{Một hình lớn được ghép từ 4 hình vuông nhỏ xếp liền nhau thành một hàng ngang. có tất cả $5$ hình vuông trong hình đó.}
{\True Một chiếc bánh sinh nhật hình tròn được cắt làm đôi bởi 1 nhát dao thẳng qua tâm. tạo thành $2$ miếng bánh.}
\loigiai{
\begin{itemize}
    \item Ý a) Sai: Kết quả đúng là $2$, không phải $3$. Số que tính dùng để xếp 2 hình vuông rời là: $4 + 4 = 8$ (que tính).\nSố que tính còn thừa là: $10 - 8 = 2$ (que tính).
\textbf{Đáp số:} $2$.
    \item Ý b) Đúng: Số que tính dùng xếp 3 hình tam giác rời là: $3 + 3 + 3 = 9$ (que tính).\nSố que tính còn thừa là: $10 - 9 = 1$ (que tính).
\textbf{Đáp số:} $1$.
    \item Ý c) Sai: Kết quả đúng là $4$, không phải $5$. Vì ghép theo hàng ngang $1 \times 4$ nên không tạo thêm hình vuông cỡ $2 \times 2$, do đó chỉ có đúng $4$ hình vuông đơn.
\textbf{Đáp số:} $4$.
    \item Ý d) Đúng: Một nhát dao thẳng cắt qua tâm chia hình tròn thành $2$ nửa hình tròn bằng nhau.
\textbf{Đáp số:} $2$.
\end{itemize}
\begin{itemize}
    \item a) \textbf{Sai}.
    \item b) \textbf{Đúng}.
    \item c) \textbf{Sai}.
    \item d) \textbf{Đúng}.
\end{itemize}
}
\end{ex}

\textbf{Phần 3. Trả lời ngắn}
\setcounterpageref{ex}{0}

\begin{ex} %[1D2V07]
Ghép 4 hình vuông nhỏ bằng nhau thành một hình vuông lớn (2 hàng, 2 cột). Hỏi có tất cả bao nhiêu hình vuông (cả hình vuông nhỏ và hình vuông lớn)?
\shortans{5}
\loigiai{
\textbf{Đáp số:} $5$\par
Có 4 hình vuông nhỏ kích thước $1 \times 1$ và 1 hình vuông lớn kích thước $2 \times 2$.\nTổng số hình vuông là: $4 + 1 = 5$ (hình vuông).
\textbf{Đáp số:} $5$.
}
\end{ex}

\begin{ex} %[1D2V07]
Bảng biển báo giao thông hình tam giác đều có 3 cạnh bằng nhau, mỗi cạnh dài 3 gang tay. Em đếm xem biển báo có bao nhiêu đỉnh nhọn?
\shortans{3}
\loigiai{
\textbf{Đáp số:} $3$\par
Biển báo hình tam giác có đúng $3$ đỉnh nhọn.
\textbf{Đáp số:} $3$.
}
\end{ex}

\begin{ex} %[1D2V07]
Mặt đồng hồ treo tường có dạng hình gì? (Hình tròn: ghi 1; Hình tam giác: ghi 2; Hình vuông: ghi 3; Hình chữ nhật: ghi 4).
\shortans{1}
\loigiai{
\textbf{Đáp số:} $1$\par
Đồng hồ mặt tròn phổ biến có dạng hình tròn, tương ứng ghi số $1$.
\textbf{Đáp số:} $1$.
}
\end{ex}

\begin{ex} %[1D2V07]
Một phong bì thư bưu điện thường có dạng hình chữ nhật. Em hãy đếm xem phong bì thư đó có bao nhiêu cạnh?
\shortans{4}
\loigiai{
\textbf{Đáp số:} $4$\par
Hình chữ nhật có $4$ cạnh.
\textbf{Đáp số:} $4$.
}
\end{ex}

\begin{ex} %[1D2V07]
Một viên gạch men lát nền nhà hình vuông có 4 cạnh. Bác thợ xếp 2 viên gạch men liền nhau tạo thành hình chữ nhật. Hỏi hình chữ nhật mới đó có bao nhiêu cạnh?
\shortans{4}
\loigiai{
\textbf{Đáp số:} $4$\par
Hình chữ nhật tạo thành cũng là một hình tứ giác có đúng $4$ cạnh ngoài bao quanh.
\textbf{Đáp số:} $4$.
}
\end{ex}

\begin{ex} %[1D2V07]
Bạn Lan có 7 hình tam giác và 3 hình vuông. Hỏi số hình tam giác nhiều hơn số hình vuông bao nhiêu hình?
\shortans{4}
\loigiai{
\textbf{Đáp số:} $4$\par
Số hình tam giác nhiều hơn số hình vuông là: $7 - 3 = 4$ (hình).
\textbf{Đáp số:} $4$.
}
\end{ex}
